\begin{lemma}[The inner domain and the continuous representative]
Let $n\ge1$, $p\ge1$, $0<\alpha\le1$, $d_0>0$, $M\ge0$ and $d>0$.  Let
$\Omega\subseteq\R^n$ be open, let $\Omega'\subseteq\R^n$ be nonempty and
bounded with
\[
d\le|x-y|\qquad\text{for every }x\in\Omega'\text{ and every }y\notin\Omega ,
\]
and let $u\in W^{1,p}(\Omega)$ with weak gradient $g$ (\noderef{MemW1pOn})
satisfy $|u|\le M$ almost everywhere on $\Omega$ and the oscillation decay
\[
\operatorname*{osc}\bigl(u,B_\rho(y)\bigr)\le4^\alpha
\max\bigl\{\operatorname*{osc}\bigl(u,B_{R_3}(y)\bigr),\ d_0R_3\bigr\}
\Bigl(\frac\rho{R_3}\Bigr)^\alpha
\qquad\bigl(\overline B_{R_3}(y)\subseteq\Omega,\ 0<\rho\le R_3\bigr)
\]
(\noderef{EssOscOn}).  Then there are an open set $\Omega_1\subseteq\R^n$ and a
function $U$ such that
\begin{itemize}
\item[(i)] $\Omega_1$ is bounded, $\overline{\Omega'}\subseteq\Omega_1$ and
  $\overline{\Omega_1}\subseteq\Omega$;
\item[(ii)] $\frac d2\le|x-y|$ for every $x\in\Omega'$ and every
  $y\notin\Omega_1$, and $\frac d2\le|x-y|$ for every $x\in\Omega_1$ and every
  $y\notin\Omega$;
\item[(iii)] $\overline B_{d/2}(x)\subseteq\Omega$ for every $x\in\Omega_1$;
\item[(iv)] $U=u$ almost everywhere on $\Omega$ and $U\in W^{1,p}(\Omega)$ with
  weak gradient $g$ (\noderef{MemW1pOn});
\item[(v)] $U$ is continuous on $\Omega_1$, and
  $U\in C^\alpha\bigl(\overline B_{d/16}(x)\bigr)$ (\noderef{MemHolderOn}) for
  every $x\in\Omega_1$.
\end{itemize}

\noindent\textbf{What this node is.}  Two things the thesis' \S4 and \S5 do
silently and that both of their main theorems need before the comparison
argument can start.  First, the passage from $\Omega'$ to an interposed inner
domain: the thesis states its Morrey and Campanato conclusions with the
subscript ``loc'' and so never names a second set, but a formalization whose
constants must depend only on the quantities the thesis lists (and in
particular on $d=d(\Omega',\partial\Omega)$ rather than on $\Omega'$ itself)
has to exhibit the intermediate set and its margins explicitly; $\Omega_1$ is
that set, and (i)--(iii) are exactly the margins the comparison argument
consumes.  Second, the passage from $u$ to a continuous representative: the
thesis evaluates $u$ at the centre of each small ball when it freezes the
operator, and it compares $v$ with $u$ on the \emph{sphere}, both of which are
meaningless for an equivalence class modulo null sets.  The De Giorgi
oscillation decay, which is the hypothesis assumed here and which the cited
H\"older estimates (C4) of the paper's Note for the formalization supply in both
settings, is what makes the choice possible, and (iv)--(v) record that
the choice costs nothing: the representative has the same weak gradient, so
every Sobolev hypothesis on $u$ transfers to it.

\noindent\textbf{On the form of the decay hypothesis.}  The hypothesis is
stated in the shape in which the cited H\"older estimates (C4) produce it ---
with the maximum against
$d_0R_3$, at an arbitrary centre, and for every radius $R_3$ whose closed ball
lies in $\Omega$ --- so that a consumer can discharge it by quoting (C4)
verbatim, and the conversion of that shape into a bound of the pure form
$K(\rho/s)^\alpha$ at the fixed scale $s=d/8$ is performed once, here.  The
hypothesis is \emph{not} that $u$ is continuous: that would beg the question
this node answers.

\noindent\textbf{On the constant.}  No constant appears in the conclusion: (v)
asserts H\"older continuity on each ball $\overline B_{d/16}(x)$ without a
norm bound, which is all the comparison argument of \S4 uses (it feeds the
H\"older datum of the Dirichlet problem, whose solution carries no constant of
the datum into the later estimates).  So no quantifier-order question arises
here, and $\alpha$ and $d_0$, which come from (C4) and therefore depend on the
consumer's $\varepsilon$, may be bound with the data.
\end{lemma}
\begin{proof}
\emph{Step 1: the two inner domains.}  Put
\[
\Omega_1=\Bigl\{x:\operatorname{dist}(x,\Omega')<\frac d2\Bigr\},
\qquad
\Omega_2=\Bigl\{x:\operatorname{dist}(x,\Omega')<\frac{3d}4\Bigr\},
\]
where $\operatorname{dist}(x,\Omega')=\inf\{|x-x'|:x'\in\Omega'\}$ is finite
because $\Omega'$ is nonempty.  Both are open, being preimages of open
half-lines under the continuous function
$x\mapsto\operatorname{dist}(x,\Omega')$, both are bounded (they are contained
in the $d$-neighbourhood of the bounded set $\Omega'$), and both contain
$\Omega'$.  For the containments and the separations we use the following
consequence of the hypothesis ``$d\le|x-y|$ for every $x\in\Omega'$ and every
$y\notin\Omega$'':
\[
0<c<d,\quad \operatorname{dist}(x,\Omega')<c,\quad y\notin\Omega
\ \Longrightarrow\ |x-y|>d-c ,
\tag{1}
\]
because such an $x$ admits $x'\in\Omega'$ with $|x-x'|<c$, and then
$d\le|x'-y|\le|x'-x|+|x-y|<c+|x-y|$.  Five consequences:
\begin{itemize}
\item $\Omega_2\subseteq\Omega$, and hence
  $\Omega'\subseteq\Omega_1\subseteq\Omega_2\subseteq\Omega$: if
  $x\in\Omega_2$ and $x\notin\Omega$ then $(1)$ with $c=\frac{3d}4$ and $y=x$
  would give $0=|x-x|>\frac d4$, which is false.
\item $\overline B_{d/8}(x)\subseteq\Omega$ for every $x\in\Omega_2$: if
  $|z-x|\le\frac d8$ then
  $\operatorname{dist}(z,\Omega')\le\operatorname{dist}(x,\Omega')+|z-x|
  <\frac{3d}4+\frac d8=\frac{7d}8$, and $(1)$ with $c=\frac{7d}8$ and $y=z$
  excludes $z\notin\Omega$ as in the previous item.
\item the separation of $\Omega_1$ from the complement of $\Omega$:
  \[
  \frac d2\le|x-y|\qquad\text{for every }x\in\Omega_1
  \text{ and every }y\notin\Omega ,
  \tag{2}
  \]
  which is $(1)$ with $c=\frac d2$ (and which $(1)$ even gives with strict
  inequality).  This is the second half of (ii).  Consequently
  $\overline B_{d/2}(x)\subseteq\Omega$ for $x\in\Omega_1$, which is (iii): a
  point $z$ with $|z-x|\le\frac d2$ and $z\notin\Omega$ would contradict the
  strict form of $(2)$ for the pair $(x,z)$.
\item $\overline{\Omega_1}\subseteq\Omega_2$ and
  $\overline{\Omega_2}\subseteq\Omega$: a point of the closure of
  $\{\operatorname{dist}(\cdot,\Omega')<c\}$ satisfies
  $\operatorname{dist}(\cdot,\Omega')\le c$, hence lies in
  $\{\operatorname{dist}(\cdot,\Omega')<c'\}$ for any $c'>c$, and the second
  containment then follows from the first item applied with $\frac{7d}8$ in
  place of $\frac{3d}4$.  Together with the first item this gives
  $\overline{\Omega_1}\subseteq\Omega$.  Also
  $\overline{\Omega'}\subseteq\Omega_1$, since every point of
  $\overline{\Omega'}$ has distance $0$ from $\Omega'$; with the boundedness of
  $\Omega_1$ and $\overline{\Omega_1}\subseteq\Omega$ this is (i).
\item if $x\in\Omega_1$ then $\overline B_{d/8}(x)\subseteq\Omega_2$, because
  $\operatorname{dist}(z,\Omega')\le\operatorname{dist}(x,\Omega')+|z-x|
  <\frac d2+\frac d8<\frac{3d}4$.
\end{itemize}
The first half of (ii) remains: if $x\in\Omega'$ and $y\notin\Omega_1$ then
$\operatorname{dist}(y,\Omega')\ge\frac d2$, and since $x\in\Omega'$ the
definition of the distance to a set gives
$|x-y|\ge\operatorname{dist}(y,\Omega')\ge\frac d2$.

\emph{Step 2: the continuous representative on $\Omega_2$.}  First,
$\operatorname*{osc}(u,B_{R_3}(y))\le2M$ whenever $B_{R_3}(y)\subseteq\Omega$:
the bound $|u|\le M$ almost everywhere on $\Omega$ restricts to $B_{R_3}(y)$,
so the essential supremum of $u$ there is at most $M$ and its essential infimum
at least $-M$ (the two facts about the essential supremum recorded in the
Preamble item on the essential supremum, applied to $u$ and to $-u$).

Put $s=\frac d8>0$ and
\[
K=4^\alpha\max\{2M,\ d_0s\}+2M\ \ge\ 2M\ \ge0 .
\]
We claim that
\[
\operatorname*{osc}\bigl(u,B_\rho(y)\bigr)\le K\Bigl(\frac\rho s\Bigr)^\alpha
\qquad\text{whenever }\overline B_\rho(y)\subseteq\Omega_2 .
\tag{3}
\]
Indeed, such a $y$ lies in $\Omega_2$, so
$\overline B_s(y)=\overline B_{d/8}(y)\subseteq\Omega$ by the second
consequence recorded in Step 1.  If $\rho\le s$, then the decay hypothesis with
$R_3=s$ --- admissible because $\overline B_s(y)\subseteq\Omega$ and
$0<\rho\le s$ --- together with the bound
$\operatorname*{osc}(u,B_s(y))\le2M$ give
$\operatorname*{osc}(u,B_\rho(y))\le4^\alpha\max\{2M,d_0s\}(\rho/s)^\alpha
\le K(\rho/s)^\alpha$.  If $\rho>s$, then
$\overline B_\rho(y)\subseteq\Omega_2\subseteq\Omega$ gives
$\operatorname*{osc}(u,B_\rho(y))\le2M\le K\le K(\rho/s)^\alpha$, because
$(\rho/s)^\alpha\ge1$ for $\rho\ge s$ and $\alpha>0$.

Now let $x_1\in\Omega_1$ and apply \noderef{OscillationDecayToHolder} with the
parameters $n,\alpha,K,M$ --- its hypotheses $1\le n$, $0<\alpha\le1$,
$0\le K$, $0\le M$ hold --- and with the data $\Omega_2$ (open), $u$, the
centre $x_1$ and the radius $s$: the containment
$\overline B_s(x_1)\subseteq\Omega_2$ is the last consequence of Step 1, the
bound $|u|\le M$ almost everywhere on $\Omega_2$ is the restriction of the
hypothesis on $\Omega$, and the decay hypothesis is exactly $(3)$.  Let
$C_2>0$ be the constant the node produces (it depends only on $n,\alpha,K,M$).
The node yields a function $U_{x_1}$ with
\begin{itemize}
\item $U_{x_1}=u$ almost everywhere on $\Omega_2$,
\item $U_{x_1}$ continuous on $\Omega_2$,
\item $U_{x_1}$ H\"older continuous of exponent $\alpha$ on
  $\overline B_{s/2}(x_1)$ (\noderef{MemHolderOn}), with
  $|U_{x_1}(x)-U_{x_1}(y)|\le C_2K\bigl(|x-y|/s\bigr)^\alpha$ there.
\end{itemize}
Any two of these functions, say $U_{x_1}$ and $U_{x_2}$ with
$x_1,x_2\in\Omega_1$, agree at every point of $\Omega_2$: their difference is
continuous on the open set $\Omega_2$ and vanishes almost everywhere there, so
the set where it does not vanish --- the preimage of $\R\setminus\{0\}$ under
a function continuous on the open set $\Omega_2$, hence an open subset of
$\Omega_2$ --- has measure zero and is therefore empty (a nonempty open set
contains a ball, whose measure is positive by the Preamble item on the volume
of a ball).  Fix now one $x_1^0\in\Omega_1$, which is possible because
$\Omega'\subseteq\Omega_1$ and $\Omega'$ is nonempty, and put
$U_0=U_{x_1^0}$; by the preceding sentence $U_{x_1}=U_0$ at every point of
$\Omega_2$, for every $x_1\in\Omega_1$.  Thus
\[
U_0=u\ \text{a.e. on }\Omega_2,
\qquad U_0\ \text{continuous on }\Omega_2,
\tag{4}
\]
and for every $x_1\in\Omega_1$,
\[
|U_0(x)-U_0(y)|\le C_2Ks^{-\alpha}|x-y|^\alpha
\qquad\bigl(x,y\in\overline B_{s/2}(x_1)\bigr),
\tag{5}
\]
the constant $C_2Ks^{-\alpha}$ being nonnegative.  Indeed
$\overline B_{s/2}(x_1)\subseteq\overline B_{d/8}(x_1)\subseteq\Omega_2$ by
the last consequence of Step 1 and $s/2=d/16<d/8$, so on
$\overline B_{s/2}(x_1)$ the function $U_0$ coincides with $U_{x_1}$, whose
H\"older bound reads
$|U_{x_1}(x)-U_{x_1}(y)|\le C_2K(|x-y|/s)^\alpha=C_2Ks^{-\alpha}|x-y|^\alpha$,
the last equality because $|x-y|\ge0$ and $s>0$.

\emph{Step 3: the modification and the transport of the Sobolev datum.}  Define
\[
U(x)=\begin{cases}U_0(x),&x\in\Omega_2,\\ u(x),&x\notin\Omega_2 .\end{cases}
\]
By $(4)$, $U=u$ almost everywhere on $\Omega$: the two functions can differ
only inside $\Omega_2$, where $U_0=u$ almost everywhere.  Apply
\noderef{AEModificationInvariance} with the parameters $n,p$ --- its hypotheses
$1\le n$, $1\le p$ hold --- on the open set $\Omega$, with $f=u$, $f'=U$ and
the field $g$, the almost-everywhere equality just established.  Its clause (2)
gives
\[
U\in W^{1,p}(\Omega)\ \text{with weak gradient }g\ (\noderef{MemW1pOn}),
\]
which together with $U=u$ almost everywhere on $\Omega$ is (iv).

Finally (v).  On $\Omega_2$ the function $U$ coincides with $U_0$ by
definition, so $U$ is continuous on $\Omega_2$ by $(4)$, hence on
$\Omega_1\subseteq\Omega_2$.  Let $x\in\Omega_1$.  By the last consequence of
Step 1, $\overline B_{d/8}(x)\subseteq\Omega_2$, and
$\overline B_{d/16}(x)=\overline B_{s/2}(x)\subseteq\overline B_{d/8}(x)$
because $s=d/8$; so on $\overline B_{d/16}(x)$ the function $U$ coincides with
$U_0$.  Therefore $U$ is continuous on $\overline B_{d/16}(x)$, and by $(5)$
applied at the centre $x$,
\[
|U(y)-U(z)|\le C_2Ks^{-\alpha}|y-z|^\alpha
\qquad\bigl(y,z\in\overline B_{d/16}(x)\bigr).
\]
With the nonnegative constant $C_2Ks^{-\alpha}$ these two facts are exactly
$U\in C^\alpha(\overline B_{d/16}(x))$ (\noderef{MemHolderOn}), which is (v).
This completes the proof.
\end{proof}
